%============================= abstract.tex================================
\chapter*{Abstract}%
\addcontentsline{toc}{chapter}{Abstract}%

Urban mass transit systems worldwide transport millions of commuters daily, necessitating extensive closed-circuit television (CCTV) surveillance networks across platform concourses, ticketing areas, and interchange corridors. However, contemporary video surveillance solutions operate predominantly as isolated, task-specific computer vision models that detect individual anomaly classes independently. Consequently, security control room operators are inundated by disjointed, unprioritized alerts that lack contextual scene comprehension, inducing severe cognitive fatigue and delaying critical emergency interventions.

To address these fundamental operational bottlenecks, this project introduces \textbf{MetroGuard AI}: a unified multi-event incident intelligence framework engineered specifically for real-time surveillance in smart metro environments. The proposed system transitions automated surveillance from siloed visual perception to context-aware spatiotemporal incident correlation and automated priority ranking. In this Project Phase-I implementation, a high-throughput, edge-capable commuter monitoring pipeline has been engineered and validated. The Phase-I architecture integrates state-of-the-art YOLOv11s object detection with ByteTrack dual-score Kalman filtering for robust commuter identity persistence across dense crowds and severe occlusions. An image-space Region of Interest (ROI) with ground-contact footpoint filtering is established to derive operational commuter occupancy levels categorized dynamically into LOW, MEDIUM, and HIGH density states. Furthermore, a Target Tracking Module (TAM) powered by MobileSAM is integrated for fine-grained silhouette segmentation and trajectory telemetry of designated individuals of interest.

The Phase-I prototype was implemented in Python using PyTorch, OpenCV, and Streamlit, and empirically benchmarked on real-world transit concourse video streams using an NVIDIA GeForce RTX 4050 Laptop GPU (CUDA-enabled). The system demonstrated superior computational efficiency, achieving an inference throughput of 62.9 FPS at 1280-pixel resolution (15.91 ms latency). Quantitative evaluation demonstrated 79.1\% mAP@0.50 on the COCO Person class with 81.2\% precision. Over an empirical benchmark of 20 sampled surveillance frames encompassing 407 ground-truth pedestrians, the model achieved a crowd counting accuracy of 99.51\% with a Weighted Absolute Percentage Error (WAPE) of only 0.49\% and a Mean Absolute Error (MAE) of 0.100 pedestrians per frame. These results confirm that the Phase-I perception backbone delivers industrial-grade precision and real-time responsiveness, providing a solid foundation for the Phase-II integration of physical altercation detection, abandoned baggage tracking, and graph-based incident fusion.

\thispagestyle{plain}
